\section{A nondegeneracy lemma from the Fargues--Fontaine curve}
\label{sec:bundles}

The formal-closure result reduces our problem to endomorphism relations
among Honda-group points. We first prove the geometric fact that will
exclude those relations: fewer than $n$ global sections of the rank-$n$,
degree-one stable bundle, independent over $\Qp$, remain independent at
the generic point. We then identify Honda-group points with sections of
these bundles. The identification holds on every perfectoid test space, so
natural maps and endomorphisms become bundle maps.

\subsection{Degree-one bundles and independent sections}

Write $X_C$ for the algebraic Fargues--Fontaine curve with coefficient
field $\Qp$. We use the vector bundle GAGA equivalence to pass between
its algebraic and adic versions. The following geometric properties
are the inputs to the determinant argument.

\begin{theorem}[Fargues--Fontaine]\label{bundles:geometry}
The curve $X_C$ is integral, regular and Noetherian of dimension one,
and $H^0(X_C,\OO_{X_C})=\Qp$.  Degrees of vector bundles are integers,
and every nonzero effective divisor has positive degree.  For each
$s\geq1$ there is a stable bundle
\[
  \mathcal E_s=\OO_{X_C}(1/s)
\]
of rank $s$ and degree one.
\end{theorem}

These assertions, including the GAGA comparison, are the
Fargues--Fontaine vector bundle theorems \cite{FF18}, in the form
of \cite[\S\,II.2.3--II.2.4, especially
Propositions~II.2.7, II.2.9--II.2.10 and Example~II.2.11]{FS21}.
The assertion about constants is
\cite[Proposition~II.2.5\textup{(ii)}]{FS21}.

\begin{lemma}[Generic independence]
\label{bundles:independence}
Let $1\leq r<n$, and let
$z_1,\ldots,z_r\in H^0(X_C,\mathcal E_n)$ be linearly independent
over $\Qp$.  Their images in the generic fiber of $\mathcal E_n$
are linearly independent over the function field $K_X$ of $X_C$.
\end{lemma}

\begin{proof}
Let $\mathcal F\subseteq\mathcal E_n$ be the saturation of the
image of the section map $\OO_{X_C}^r\to\mathcal E_n$, and let
$b=\rank\mathcal F$.  On a regular Noetherian curve saturation is
a subbundle, and $1\leq b\leq r<n$.  Choose $b$ of the given
sections that are independent at the generic point.  Their wedge is
a nonzero global section of $\det\mathcal F$, so
$\deg\mathcal F\geq0$.  Semistability of $\mathcal E_n$ gives
\[
  0\leq\deg\mathcal F\leq \frac{b}{n}<1.
\]
Its integral degree is therefore zero.  The divisor of the wedge
section is effective of degree zero, hence empty by
Theorem~\ref{bundles:geometry}.  The chosen $b$ sections
thus trivialize $\mathcal F$ by the determinant criterion.

All the original sections belong to $H^0(X_C,\mathcal F)$.  Under
this trivialization that space is $\Qp^b$, so their assumed
$\Qp$-independence gives $r\leq b$.  Consequently $b=r$, as
required.
\end{proof}

To apply this lemma to points of $\Gamma_s(\mm_C)$, we need their
realization as global sections and a corresponding realization of
natural maps as bundle maps. We establish both comparisons next.

\subsection{Honda groups as relative section sheaves}

We use the big pro-\'etale site of perfectoid spaces over
$C=\Cp^\flat$, denoted by $\Perf_C$.  For an affinoid
$S=\operatorname{Spa}(R,R^+)$, let $R^{\circ\circ}$ be the
topologically nilpotent elements of $R$, and put
\[
  \mathcal H_s(S)=\Gamma_s(R^{\circ\circ}),\qquad s\geq 1.
\]
The underlying functor is the perfected open unit disc, equipped with the
Honda group law; in particular, it is a sheaf.  Since $R$ is perfect,
$[p]_{\Gamma_s}(z)=z^{q_s}$ is bijective.  Together with the continuous
formal $\Zp$-action, this makes $\mathcal H_s$ a sheaf of
$\Qp$-vector spaces.  All roots appearing in this description are unique.

The choice of $R^+$ causes no restriction on these points.  Indeed,
if $z\in R^{\circ\circ}$, openness of $R^+$ gives $z^m\in R^+$
for some $m\geq1$.  The element $z$ is integral over $R^+$ and therefore
belongs to $R^+$, which is integrally closed in $R$.  Thus
$R^{\circ\circ}\subseteq R^+$ for every affinoid perfectoid pair.
We also write $\mathcal H_s(R)$ for these values, and
$\mathcal H_s(C)=\Gamma_s(\mm_C)$.

For $S\in\Perf_C$, write $X_S$ for the relative curve.
If $\mathcal V$ is a vector bundle on $X_C$, let
$\mathcal V_S$ be its pullback and define
\[
  \mathcal B(\mathcal V)(S)=H^0(X_S,\mathcal V_S).
\]
This is a pro-\'etale sheaf; in fact relative sections satisfy
$v$-descent \cite[Proposition~II.2.1]{FS21}.

The section-sheaf dictionary has two parts: it identifies the Honda
sheaves with the bundles just introduced, and it recovers bundle maps
from additive maps of their section sheaves.

\begin{theorem}[Fargues--Fontaine and Le Bras]
\label{bundles:input}
\leavevmode
\begin{enumerate}
\item[\textup{(i)}]
For every $s\geq1$, after choosing an identification of a height-$s$
Lubin--Tate special fiber with $\Gamma_s$, there is a natural isomorphism of
$\Qp$-vector space sheaves
\[
  \mathcal H_s\ \cong\ \mathcal B(\mathcal E_s).
\]
In particular, its value at $\operatorname{Spa}(C,\OO_C)$ identifies
$\Gamma_s(\mm_C)$ with $H^0(X_C,\mathcal E_s)$.

\item[\textup{(ii)}]
For vector bundles $\mathcal V,\mathcal W$ on $X_C$ with nonnegative
Harder--Narasimhan slopes, the natural map
\[
  \Hom_{X_C}(\mathcal V,\mathcal W)
  \longrightarrow
  \Hom_{\operatorname{Sh}_{\mathrm{Ab}}(\Perf_C)}
       \bigl(\mathcal B(\mathcal V),\mathcal B(\mathcal W)\bigr)
\]
is an isomorphism.  Thus every additive sheaf morphism between these
section sheaves is induced by a unique bundle morphism.
\end{enumerate}
\end{theorem}

\begin{proof}[The identification and full faithfulness]
For part~\textup{(i)}, let $E^{(s)}\subset\Cp$ be the unramified
extension of $\Qp$ of degree $s$, and let $G^{(s)}$ be a Lubin--Tate
formal group for $E^{(s)}$, with uniformizer $p$
\cite[Theorem~1]{LT65}.  Its reduction over
$k$ has height $s$ as a $p$-divisible formal group.  Choose a formal
group isomorphism from this reduction to $\Gamma_s$.

Tilting gives an equivalence between the big pro-\'etale sites over
$\Cp$ and over $C$ \cite[Theorem~2.7]{LeBras18}.  Thus each
$S=\operatorname{Spa}(R,R^+)$ has its specified untilt
$S^\sharp=\operatorname{Spa}(R^\sharp,R^{\sharp+})$ over $\Cp$.
Fargues--Scholze identify the Lubin--Tate universal cover with relative
sections by a natural isomorphism
\cite[Proposition~II.2.2]{FS21}
\[
  \varprojlim_{[p]}
     G^{(s)}\bigl((R^\sharp)^{\circ\circ}\bigr)
  \ \cong\
  H^0\bigl(X_{S,E^{(s)}},\OO(1)\bigr).
\]
Here $X_{S,E^{(s)}}$ denotes the relative curve with coefficient field
$E^{(s)}$.  Their integral formulation gives the same universal-cover
points, since the topologically nilpotent elements belong to every
allowed integral subring.

We spell out the comparison of the left side with $\mathcal H_s(S)$.
Choose $p^\flat\in\OO_C$ with $(p^\flat)^\sharp=p$.
The tilting comparison identifies
\[
  (R^\sharp)^\circ/p
    \ \cong\ R^\circ/(p^\flat),
\]
with compatible residue embeddings of $k$.  On these quotient rings
the chosen special-fiber isomorphism identifies the two formal group
laws.  It remains to pass between universal-cover points before and
after reduction.

The relevant reduction comparison can be seen directly.  Suppose an
integral formal group law is evaluated in a complete uniform algebra,
and reduce its power-bounded ring modulo a fixed small base parameter
of norm $\delta<1$.  A nilpotent element of the quotient lifts to a
topologically nilpotent element: if a lift $a$ satisfies
$a^m$ in the parameter ideal, its spectral norm satisfies
$\|a\|^m\leq\delta$.  Given compatible formal points
$(\bar a_j)_{j\geq0}$ of the quotient under $[p]$, choose such lifts
$a_j$.  For each fixed $i$, the sequence
\[
  [p]^j(a_{i+j}),\qquad j\geq0,
\]
converges.  Indeed the formal difference between two successive terms
is obtained by pushing down $j$ times a difference of norm at most
$\delta$.  For a difference of norm $\epsilon<1$, integrality and
the linear coefficient $p$ give the bound
\[
  \|[p](w)\|
     \leq\max\{|p|\epsilon,\epsilon^2\}
     \qquad (\|w\|\leq\epsilon),
\]
and its iterates tend to zero.  In characteristic $p$ the linear term
is zero, giving the same conclusion.  These estimates use an
equivalent complete spectral norm, which is submultiplicative and
power-multiplicative; no field hypothesis on the test algebra is
needed.  Integral formal translation also turns these bounds on
formal differences into bounds on ordinary coordinate differences.

For fixed $i$ the limit has norm at most
$\max\{\|a_i\|,\delta\}<1$.  The scalar ideal is closed for the
spectral norm, so the limits retain the prescribed reductions.
The limits therefore form a compatible
lift of $(\bar a_j)$ by topologically nilpotent points.  Two choices of
lifts give the same limits by the identical contraction estimate, and
the same estimate proves uniqueness of a compatible lift.  The
comparison therefore commutes with algebra maps and with the group
operations.  It also commutes with formal scalar multiplication.
This proves the required natural reduction invariance of the
universal cover; compare \cite[\S\,II.2.1]{FS21} and
\cite[Proposition~2.14, Lemma~2.15, and Remark~2.21]{LeBras18}.

Applying this comparison on both sides of the displayed tilting
isomorphism gives
\[
  \varprojlim_{[p]}
     G^{(s)}\bigl((R^\sharp)^{\circ\circ}\bigr)
  \ \cong\
  \varprojlim_{[p]}\Gamma_s(R^{\circ\circ})
  \ \cong\ \Gamma_s(R^{\circ\circ}).
\]
The last isomorphism sends $z$ to the compatible sequence
$(z^{1/q_s^j})_{j\geq0}$ and uses perfectness of $R$.

We now pass from the coefficient field $E^{(s)}$ to $\Qp$.
Along the unramified coefficient extension
$\pi_s:X_{S,E^{(s)}}\to X_S$, one has
\[
  (\pi_s)_*\OO(1)\cong(\mathcal E_s)_S.
\]
This is the unramified pushforward description of
$\OO(1/s)$ \cite[\S\,5.1]{LeBras18}; see also
\cite[proofs of Theorem~II.2.14 and Corollary~II.3.3]{FS21}.
The description is relative and natural
in $S$: on the Frobenius cover, pushforward cycles $s$ lines, with
the $s$-fold isocrystal Frobenius having multiplier $p^{-1}$.
It therefore gives the rank-$s$ isocrystal of slope $-1/s$, and hence
the bundle of slope $+1/s$.  Taking sections commutes with this
pushforward, proving part~\textup{(i)} on the whole site.

For part~\textup{(ii)}, Le Bras identifies Banach--Colmez sheaves with
the tilted heart
\[
  \operatorname{Coh}_{X_C}^{-}
  =\left\{F\in D^b(\operatorname{Coh}_{X_C}):
       H^i(F)=0\ (i\ne-1,0),\quad
       H^{-1}(F)<0,\quad H^0(F)\geq0\right\}
\]
by the functor $R^0\tau_*$
\cite[Theorem~1.2, Theorem~7.1]{LeBras18}.  The inequalities here
refer to all Harder--Narasimhan slopes.  A nonnegative-slope vector
bundle, placed in degree zero, belongs to this heart, and its
$R^0\tau_*$ is the relative section sheaf just defined
\cite[Proposition~6.8 and its proof]{LeBras18}.
The heart is full in the derived category, so Hom between two such
bundles is their usual bundle Hom.  Moreover, the proof of
\cite[Theorem~7.1]{LeBras18} establishes the degree-zero Ext
comparison in the category of abelian sheaves on the big site.
Consequently it applies to the Hom group in part~\textup{(ii)},
without an additional analytic-morphism or scalar-continuity
hypothesis.  The fixed tilting equivalence transports this statement
to $\Perf_C$.
\end{proof}

In particular, every endomorphism in $A_s$ acts naturally on
$\mathcal H_s$ and therefore on $\mathcal E_s$.  This action extends
to $\Delta_s=A_s[1/p]$, since multiplication by $p$ is invertible on
both sides.  We will use this action on the bundle itself.  No full
faithfulness assertion for the isocrystal-to-bundle functor is needed.

\subsection{Endomorphism orbits}

\begin{proposition}[Nondegeneracy]
\label{bundles:nondegeneracy}
Let $0<h<n$ and $1\leq r<n$.  Suppose that
$z_1,\ldots,z_r\in\mathcal H_n(C)$ are $\Qp$-linearly independent
and that
\[
  \nu:\mathcal H_n\longrightarrow\mathcal H_h
\]
is a nonzero morphism of $\Qp$-vector space sheaves on $\Perf_C$.
Then the subgroup
\[
  S_\nu=\bigl\{\bigl(\nu(dz_1),\ldots,\nu(dz_r)\bigr):
                    d\in A_n\bigr\}
       \subseteq\Gamma_h^r(\mm_C)
\]
satisfies no nonzero formal power series equation over $k$.
\end{proposition}

\begin{proof}
The set $S_\nu$ is a subgroup because the ring addition in $A_n$
and the map $\nu$ respect the formal group laws.  If a nonzero
$k$-formal series vanished on it,
Proposition~\ref{formal:closure} would give $a_1,\ldots,a_r\in A_h$,
not all zero, such that
\begin{equation}
\label{bundles:annihilator-relation}
  \sum_{i=1}^r a_i\nu d(z_i)=0
          \qquad\text{for every }d\in A_n.
\end{equation}
The sum in this equation is the target group sum.

By Theorem~\ref{bundles:input}, regard the $z_i$ as global sections
of $\mathcal E_n$ and $\nu$ as a nonzero bundle map
$\mathcal E_n\to\mathcal E_h$.  The operators in $A_n$ and $A_h$
also act by bundle maps.  Passing to the generic point gives a
unital $K_X$-algebra homomorphism
\[
  \Delta_n\otimes_{\Qp}K_X
      \longrightarrow
  \End_{K_X}\bigl((\mathcal E_n)_{K_X}\bigr).
\]
The source is central simple over $K_X$.  The homomorphism is
therefore injective, and both algebras have dimension $n^2$ over
$K_X$, so it is an isomorphism.  In particular, splitting of
$\Delta_n$ over $K_X$ follows from the displayed action; it has
not been assumed in constructing that action.

Equation~\eqref{bundles:annihilator-relation} is linear in $d$.
First extend it from $A_n$ to $A_n[1/p]=\Delta_n$.  At the generic
point it then extends $K_X$-linearly to all of the displayed matrix
algebra.  This is the extension of a vanishing linear identity; no
assertion that an arbitrary matrix preserves the order $A_n$ is
involved.

By Lemma~\ref{bundles:independence}, the generic vectors $z_i$ are
independent.  Choose an index $i_0$ with $a_{i_0}\ne0$.  This
operator is invertible in the division algebra $\Delta_h$, hence
acts invertibly on $\mathcal E_h$.  The generic map $\nu_{K_X}$
is nonzero: a nonzero morphism between vector bundles on an integral
curve cannot have zero generic fiber.  Choose
$w\in(\mathcal E_n)_{K_X}$ with $\nu_{K_X}(w)\ne0$, and choose a
matrix $M$ satisfying
\[
  Mz_{i_0}=w,\qquad Mz_i=0\quad(i\ne i_0).
\]
Such a matrix exists by independence.  Substituting it into the
extended identity gives
$a_{i_0}\nu_{K_X}(w)=0$, a contradiction.
\end{proof}

The hypothesis that $\nu$ is a sheaf morphism is important: the
bundle argument does not apply to an arbitrary linear map between
the $C$-valued groups.  In the next section we construct the required
nonzero natural morphism from a deformation point, together with
the independent vectors to which
Proposition~\ref{bundles:nondegeneracy} will be applied.
